\documentclass[12pt, letterpaper]{article}

\usepackage[utf8]{inputenc}
\usepackage[T1]{fontenc}
\usepackage[spanish]{babel}
\usepackage{geometry}
\usepackage{amsmath, amssymb, amsthm}
\usepackage{booktabs}
\usepackage{array}
\usepackage{microtype}
\usepackage{setspace}
\usepackage{parskip}
\usepackage{titlesec}
\usepackage{fancyhdr}
\usepackage{enumitem}
\usepackage{bm}
\usepackage{mathtools}
\usepackage{tcolorbox}
\usepackage{hyperref}

\tcbuselibrary{skins, breakable}

\geometry{top=2.8cm, bottom=3.0cm, left=3.2cm, right=3.2cm}
\setstretch{1.20}

\pagestyle{fancy}
\fancyhf{}
\fancyhead[C]{\small\textsc{Econometria · Gujarati --- Ejercicio 10.11}}
\fancyfoot[C]{\small\thepage}
\renewcommand{\headrulewidth}{0.4pt}
\renewcommand{\footrulewidth}{0pt}

\titleformat{\section}
  {\large\bfseries\scshape}{\thesection.}{0.6em}{}
  [\vspace{-4pt}\rule{\linewidth}{0.4pt}]
\titleformat{\subsection}{\normalsize\bfseries}{\thesubsection.}{0.5em}{}
\titleformat{\subsubsection}{\normalsize\itshape}{}{0em}{}

\newtheorem{prop}{Proposicion}
\newtheorem{lema}{Lema}
\theoremstyle{remark}
\newtheorem{obs}{Observacion}

\newtcolorbox{clave}{
  colback=white, colframe=black,
  boxrule=0.5pt, arc=3pt,
  left=8pt, right=8pt, top=6pt, bottom=6pt,
  breakable
}

\newcommand{\E}{\mathbb{E}}
\newcommand{\Var}{\operatorname{Var}}
\newcommand{\Cov}{\operatorname{Cov}}
\newcommand{\plim}{\operatorname{plim}}

\begin{document}

\begin{center}
  {\LARGE\bfseries Regresion por Pasos:}\\[6pt]
  {\LARGE\bfseries Seleccion de Variables, Multicolinealidad}\\[6pt]
  {\LARGE\bfseries y sus Problemas Formales}\\[12pt]
  {\large\itshape Por que los procedimientos mecanicos de seleccion son inadecuados}\\[14pt]
  \rule{0.55\linewidth}{0.6pt}\\[8pt]
  {\normalsize Ejercicio~10.11 --- \textit{Econometria}, Gujarati \& Porter, 5.a ed.}\\[4pt]
  {\small\textsc{Analisis formal: sesgo por pretesting, no monotonía de la SCE,
  inflacion del error tipo I y alternativas validas}}
\end{center}

\vspace{10pt}

\noindent\textbf{Planteamiento.}\quad
La \textbf{regresion por pasos} (\textit{stepwise regression}) es un
procedimiento automatico de seleccion de variables que opera de dos formas:

\begin{itemize}[leftmargin=1.6em, itemsep=4pt]
  \item \textbf{Hacia delante} (\textit{forward}): se parte del modelo vacio
    y se agrega una variable a la vez, eligiendo en cada paso la que mas
    incrementa la SCE, siempre que el estadistico $F$ correspondiente supere
    un umbral predefinido.
  \item \textbf{Hacia atras} (\textit{backward}): se parte del modelo completo
    con todas las variables disponibles y se elimina una a la vez, comenzando
    por la de menor contribucion marginal a la SCE, hasta que todas las
    variables restantes sean significativas segun la prueba $F$.
\end{itemize}

\noindent Se pregunta si estos procedimientos son recomendables a la luz
de lo que se sabe sobre multicolinealidad.

\noindent\textbf{Respuesta directa: no se recomiendan}, por razones formales
que se desarrollan a continuacion.

\section{La logica de la regresion por pasos: criterio de la SCE}

\subsection{El estadistico $F$ incremental}

En cada paso del procedimiento hacia delante, se evalua si agregar la
variable $X_j$ mejora significativamente el ajuste. El estadistico $F$
incremental compara el modelo sin $X_j$ (restringido, con $R^2_R$) contra
el modelo con $X_j$ (no restringido, con $R^2_{NR}$):

\begin{equation}
  F^* = \frac{(R^2_{NR} - R^2_R)/1}{(1 - R^2_{NR})/(n - k)}
  \label{eq:F_increm}
\end{equation}

Si $F^* > F_c$ (valor critico al nivel $\alpha$), $X_j$ se incluye. Este
estadistico es equivalente al cuadrado del estadistico $t$ individual de
$X_j$ en el modelo que ya incluye las demas variables. El procedimiento es
entonces secuencial: cada decision depende del conjunto de variables que
ya estan en el modelo.

\subsection{La SCE como criterio de seleccion}

La suma de cuadrados explicada (SCE) del modelo con $k$ regresores es:

\begin{equation}
  \text{SCE}_k = \sum_{i=1}^n (\hat Y_i - \bar Y)^2
  \label{eq:SCE}
\end{equation}

Agregar cualquier variable al modelo nunca puede reducir la SCE, porque
MCO minimiza la suma de cuadrados residual y la SCE es su complemento.
Por tanto:

\begin{equation}
  \text{SCE}_{k+1} \geq \text{SCE}_k \quad \text{siempre}
  \label{eq:monotonia_SCE}
\end{equation}

El procedimiento por pasos elige las variables que maximizan el incremento
marginal en la SCE en cada paso. El problema es que bajo multicolinealidad,
este criterio depende criticamente del orden de entrada, como se demuestra
a continuacion.

\section{Problema 1: dependencia del orden y arbitrariedad}

\subsection{Demostracion de la no unicidad del resultado}

Sea un modelo con tres variables candidatas: $X_2$, $X_3$ y $X_4$, donde
$r_{23}$ es alta. La contribucion marginal de $X_3$ al agregar al modelo
que ya contiene $X_2$ es:

\begin{equation}
  \Delta\text{SCE}(X_3 | X_2)
    = \text{SCE}(X_2, X_3) - \text{SCE}(X_2)
  \label{eq:delta_SCE_cond}
\end{equation}

Pero si en cambio se agrega $X_3$ primero y luego se evalua $X_2$:

\begin{equation}
  \Delta\text{SCE}(X_2 | X_3)
    = \text{SCE}(X_2, X_3) - \text{SCE}(X_3)
  \label{eq:delta_SCE_cond2}
\end{equation}

En general, $\Delta\text{SCE}(X_3 | X_2) \neq \Delta\text{SCE}(X_3)$
porque el efecto marginal de $X_3$ cambia segun que otras variables esten
ya en el modelo. Con colinealidad alta ($r_{23} \approx 1$), la contribucion
marginal de $X_3$ dado $X_2$ es casi cero (ambas explican casi lo mismo),
mientras que la contribucion marginal de $X_3$ sin $X_2$ puede ser grande.

\subsection{Implicacion para la seleccion de modelos}

Suponga que se tienen dos posibles modelos equivalentes en ajuste:

\begin{align}
  \text{Modelo A:} \quad Y &= \beta_1 + \beta_2 X_2 + u \\
  \text{Modelo B:} \quad Y &= \gamma_1 + \gamma_3 X_3 + v
\end{align}

con $r_{23} = 0.98$ y $R^2_A \approx R^2_B$. El procedimiento hacia delante
elegira uno u otro dependiendo unicamente de cual explique marginalmente
un poco mas en el primer paso, aunque ambos sean virtualmente equivalentes
en ajuste y en contenido informativo. La eleccion es arbitraria y no tiene
ningun respaldo teorico.

\begin{clave}
  Bajo multicolinealidad, la contribucion marginal de una variable a la
  SCE depende criticamente del conjunto de variables ya incluidas. El resultado
  del procedimiento por pasos varia con el orden de entrada o salida,
  produciendo selecciones arbitrarias de modelos.
\end{clave}

\section{Problema 2: inflacion del error tipo I por pruebas multiples}

\subsection{El problema del pretesting}

Cuando se realizan $m$ pruebas de hipotesis secuenciales sobre los mismos
datos, el nivel de significancia nominal $\alpha$ de cada prueba individual
ya no refleja la probabilidad real de cometer un error tipo~I en el proceso
total. Si cada prueba tiene nivel $\alpha$, la probabilidad de cometer al
menos un error tipo~I en $m$ pruebas independientes es:

\begin{equation}
  P(\text{al menos un error tipo I}) = 1 - (1-\alpha)^m
  \label{eq:bonferroni}
\end{equation}

Para $\alpha = 0.05$ y $m = 10$ variables candidatas:

\begin{equation}
  1 - (1-0.05)^{10} = 1 - 0.5987 \approx 0.40
  \label{eq:tipo1_10}
\end{equation}

Con nivel nominal del $5\%$ por prueba, la probabilidad real de incluir al
menos una variable irrelevante en 10 pasos es del $40\%$.

\subsection{Distribucion del estadistico bajo seleccion}

El estadistico $t$ del $j$-esimo coeficiente en el modelo final seleccionado
por el procedimiento por pasos \textbf{no sigue la distribucion $t$ de Student}
con los grados de libertad habituales. La razon es que el modelo fue elegido
precisamente porque los estadisticos $t$ superaron un umbral: los $t$ del
modelo seleccionado estan condicionados a haber sido seleccionados, lo que
sesga su distribucion hacia valores grandes:

\begin{equation}
  \E[|t_j| \,|\, \text{variable seleccionada}] > \E[|t_j|]
  \label{eq:t_sesgado}
\end{equation}

Esto invalida la inferencia estandar: los errores estandar, intervalos de
confianza y p-valores del modelo final son anticonservativos.

\subsection{El problema del data mining}

Buscar el ``mejor'' modelo entre todos los subconjuntos posibles de $k$
variables candidatas implica evaluar hasta $2^k - 1$ modelos. Para $k = 10$:
$2^{10} - 1 = 1023$ modelos. Cada uno es una prueba implicita de hipotesis.
La probabilidad de encontrar un modelo con $R^2$ alto por pura casualidad
—incluso si todas las variables son irrelevantes— es alta. Este fenomeno
se denomina \textbf{sobreajuste por seleccion} (\textit{selection overfitting}).

\begin{clave}
  Los procedimientos por pasos realizan multiples pruebas de hipotesis sobre
  los mismos datos. El nivel nominal $\alpha$ no refleja la probabilidad real
  de error tipo~I, que puede superar el $40\%$ con 10 variables candidatas.
  La inferencia del modelo seleccionado es invalida.
\end{clave}

\section{Problema 3: sesgo por variable omitida bajo multicolinealidad}

\subsection{Mecanismo del sesgo}

Si una variable teoricamente relevante $X_j$ esta altamente correlacionada
con otras variables ya en el modelo, su contribucion marginal a la SCE es
pequena, pues las otras variables ya ``explican'' gran parte de su variacion.
En el procedimiento hacia delante, $X_j$ puede no ser seleccionada aunque
$\beta_j \neq 0$. En el procedimiento hacia atras, puede ser eliminada
porque su estadistico $t$ no es significativo (inflado por colinealidad).

La formula del sesgo por variable omitida (derivada en ejercicios anteriores)
se aplica aqui de forma directa:

\begin{equation}
  \E[\hat\beta_{\text{red},j}] = \beta_j + \sum_{k \notin \mathcal{M}} \beta_k b_{jk}
  \label{eq:sesgo_omision}
\end{equation}

donde $\mathcal{M}$ es el conjunto de variables omitidas y $b_{jk}$ son los
coeficientes de las regresiones auxiliares de las variables omitidas sobre
las incluidas. Bajo colinealidad, $b_{jk}$ es grande, amplificando el sesgo.

\subsection{Paradoja: eliminacion de variables porque su $t$ es pequeno por colinealidad}

La colinealidad infla los errores estandar de todos los regresores
correlacionados, reduciendo sus estadisticos $t$:

\begin{equation}
  t_j = \frac{\hat\beta_j}{\text{ee}(\hat\beta_j)}
       = \frac{\hat\beta_j}{\hat\sigma/\sqrt{S_{jj}}} \cdot \frac{1}{\sqrt{\text{VIF}_j}}
  \label{eq:t_VIF}
\end{equation}

Con VIF$_j$ alto, $t_j$ puede ser pequeno aunque $\beta_j \neq 0$. El
procedimiento por pasos interpretaria esto como evidencia de que $X_j$ es
irrelevante y la eliminaria, cuando en realidad es la colinealidad la que
impide detectar su efecto. El resultado es la omision de variables relevantes
y el sesgo de sus coeficientes en las variables que permanecen.

\begin{clave}
  En presencia de multicolinealidad, los procedimientos por pasos tienden a
  eliminar variables teoricamente relevantes porque su efecto individual no
  es detectado por la prueba $F$ o $t$, introduciendo sesgo por variable omitida.
\end{clave}

\section{Problema 4: inestabilidad de los coeficientes}

Cuando la multicolinealidad es severa, el estimador MCO es numericamente
inestable: pequenas perturbaciones en los datos producen cambios grandes
en los coeficientes. La adicion o eliminacion de una variable en el
procedimiento por pasos puede desencadenar cambios drasticos en todos los
demas coeficientes del modelo:

\begin{equation}
  \frac{\partial \hat\beta_j}{\partial X_{ki}}
  \propto [(\mathbf{X'X})^{-1}]_{jk}
  \label{eq:sensibilidad}
\end{equation}

Con $\det(\mathbf{X'X}) \approx 0$, los elementos de $(\mathbf{X'X})^{-1}$
son grandes y la sensibilidad es extrema. El modelo seleccionado en una
muestra puede ser muy diferente del que se seleccionaria en una muestra
ligeramente diferente del mismo fenomeno.

\section{Alternativas recomendadas}

Dado que los procedimientos mecanicos son inadecuados, se presentan las
alternativas metodologicamente correctas.

\subsection{Especificacion teorica del modelo}

La especificacion debe basarse en la teoria economica, no en criterios
estadisticos. Un modelo con $k$ variables relevantes es preferible a uno
con $k' < k$ variables aunque el segundo tenga menor varianza de estimacion,
porque el primero es insesgado y el segundo no.

\subsection{Criterios de informacion penalizados}

Los criterios de informacion penalizan la complejidad del modelo y son
apropiados para comparar modelos con diferente numero de parametros:

\begin{align}
  \text{AIC} &= -2\ell(\hat{\boldsymbol\beta}) + 2k
  \label{eq:AIC}\\
  \text{BIC} &= -2\ell(\hat{\boldsymbol\beta}) + k\ln n
  \label{eq:BIC}
\end{align}

donde $\ell$ es el logaritmo de la verosimilitud maximizada y $k$ es el
numero de parametros. El BIC penaliza mas que el AIC y tiende a seleccionar
modelos mas parsimoniosos. A diferencia del procedimiento por pasos, estos
criterios comparan modelos \emph{globalmente} sin sesgo de pretesting.

\subsection{Validacion cruzada}

Para evitar el sobreajuste, se puede usar validacion cruzada (\textit{cross-validation}):
dividir la muestra en subconjuntos de entrenamiento y prueba, estimar el
modelo en el subconjunto de entrenamiento y evaluar su ajuste en el de prueba.
El modelo que minimiza el error de prediccion fuera de muestra es preferido.

\subsection{Regularizacion: regresion Ridge y LASSO}

Los metodos de regularizacion encogen los coeficientes hacia cero de forma
controlada, resolviendo simultaneamente el problema de multicolinealidad y
la seleccion de variables:

\begin{align}
  \hat{\boldsymbol\beta}^{\text{Ridge}} &= \arg\min_{\boldsymbol\beta}
    \left\{ \sum_i (Y_i - \mathbf{x}_i'\boldsymbol\beta)^2
            + \lambda \sum_j \beta_j^2 \right\}
  \label{eq:Ridge}\\
  \hat{\boldsymbol\beta}^{\text{LASSO}} &= \arg\min_{\boldsymbol\beta}
    \left\{ \sum_i (Y_i - \mathbf{x}_i'\boldsymbol\beta)^2
            + \lambda \sum_j |\beta_j| \right\}
  \label{eq:LASSO}
\end{align}

El LASSO produce soluciones con coeficientes exactamente iguales a cero
para variables irrelevantes (seleccion automatica) pero con fundamento
estadistico riguroso, sin los problemas de pretesting del procedimiento
por pasos.

\section{Sintesis: cuadro comparativo}

\begin{table}[ht]
\centering
\caption{Comparacion de metodos de seleccion de variables}
\label{tab:comparacion}
\small
\begin{tabular}{p{3.5cm} p{3.5cm} p{3.5cm} p{3cm}}
\toprule
\textbf{Criterio}
  & \textbf{Regresion por pasos}
  & \textbf{AIC/BIC}
  & \textbf{LASSO} \\
\midrule
Sesgo por pretesting
  & Alto (pruebas multiples)
  & Bajo (comparacion global)
  & Ninguno \\[4pt]
Efecto de la colinealidad
  & Elimina variables relevantes
  & Penaliza complejidad; no resuelve colinealidad
  & Encoge coeficientes; maneja colinealidad \\[4pt]
Unicidad del resultado
  & No: depende del orden
  & Si: criterio unico
  & Si: para $\lambda$ fijo \\[4pt]
Validez de inferencia
  & Invalida (distribucion $t$ sesgada)
  & Valida si se escoge correctamente
  & Requiere inferencia post-seleccion \\[4pt]
Base teorica
  & Ninguna: puramente mecanico
  & Teoria de la informacion
  & Teoria estadistica de regularizacion \\
\bottomrule
\end{tabular}
\end{table}

\begin{clave}
  \textbf{Conclusion.} La regresion por pasos \textbf{no se recomienda} en
  econometria aplicada, especialmente bajo multicolinealidad, por cuatro
  razones formales: (1) el resultado depende arbitrariamente del orden de
  entrada o salida; (2) la realizacion de multiples pruebas de hipotesis
  secuenciales infla el error tipo~I hasta valores del 40\% o mas; (3) las
  variables relevantes pueden ser eliminadas porque la colinealidad reduce
  sus estadisticos $t$, introduciendo sesgo por variable omitida; (4) los
  coeficientes del modelo seleccionado son inestables y la inferencia
  estandar es invalida. Las alternativas correctas son: especificacion
  teorica previa, criterios AIC/BIC, validacion cruzada y metodos de
  regularizacion como Ridge o LASSO.
\end{clave}

\appendix

\section{Demostracion de la inflacion del error tipo I bajo seleccion secuencial}

Sean $m$ variables candidatas independientes (para simplificar). La hipotesis
nula para cada una es $H_{0j}: \beta_j = 0$. La probabilidad de no cometer
un error tipo~I en ninguna de las $m$ pruebas es:

\begin{equation}
  P(\text{ningun error tipo I}) = (1-\alpha)^m
  \label{eq:no_error}
\end{equation}

Por tanto, la probabilidad de cometer al menos un error tipo~I es:

\begin{equation}
  \alpha_{\text{exp}} = 1 - (1-\alpha)^m
  \label{eq:alpha_exp}
\end{equation}

Esta es la \textbf{tasa de error familiar} (\textit{familywise error rate},
FWER). Para controlarla al nivel $\alpha$ se debe usar la correccion de
Bonferroni: realizar cada prueba al nivel $\alpha/m$, de modo que:

\begin{equation}
  \alpha_{\text{exp}} \leq 1 - \left(1 - \frac{\alpha}{m}\right)^m \approx \alpha
  \label{eq:bonferroni_corr}
\end{equation}

El procedimiento por pasos no aplica ninguna correccion de este tipo, lo que
hace que su FWER sea sistematicamente superior al nivel nominal. Los p-valores
reportados en el modelo final deben interpretarse con extrema cautela.

\section{No monotonia de las contribuciones marginales bajo colinealidad}

Sea el modelo verdadero $Y = \beta_2 X_2 + \beta_3 X_3 + u$ con $r_{23} = r$.
La contribucion marginal de $X_3$ dado $X_2$ es:

\begin{equation}
  \Delta R^2(X_3|X_2)
    = \frac{r_{Y,X_3\cdot X_2}^2 (1 - R^2_{Y,X_2})}{1}
  \label{eq:delta_R2}
\end{equation}

donde $r_{Y,X_3 \cdot X_2}$ es la correlacion parcial de $Y$ con $X_3$
dado $X_2$. Esta correlacion parcial puede ser muy diferente de la correlacion
simple $r_{Y,X_3}$ cuando $r_{23}$ es alta, incluso si $\beta_3 \neq 0$.
Con $r_{23} \approx 1$, la correlacion parcial tiende a cero y $X_3$ no sera
seleccionada en el segundo paso, aunque sea relevante. Esta es la manifestacion
formal de por que el orden importa.

\vspace{14pt}
\noindent\rule{\linewidth}{0.4pt}

\noindent{\small\textit{Nota.} La critica a los metodos de regresion por pasos
esta bien documentada en la literatura econometrica y estadistica. Hocking (1976),
Derksen y Keselman (1992) y Harrell (2001) documentan empiricamente las altas
tasas de inclusion de variables irrelevantes y exclusion de variables relevantes
bajo estos procedimientos. Para seleccion formal de variables con control del
error tipo~I, se recomienda el procedimiento de Benjamini-Hochberg (1995) para
controlar la tasa de falsos descubrimientos (\textit{False Discovery Rate}).}

\end{document}
